\begin{figure}[htbp]
\centering
\begin{tikzpicture}[x=1.45cm,y=1.45cm,font=\small,>=Latex]
 \draw[black!55] (-2.1,0)--(2.65,0);
 \draw[dashed,black!50] (0,-.12)--(0,3.05);
 \node[above] at (0,3.05) {cut};
 \node at (-1.3,2.95) {left source state};
 \node at (1.65,2.95) {right source state};
 \draw[very thick,blue!65!black]
   (-1.65,0) .. controls (-1.65,.38) and (-.45,.45) .. (0,.45)
   .. controls (.58,.45) and (.58,.9) .. (0,.9)
   .. controls (-.58,.9) and (-.58,1.35) .. (0,1.35)
   .. controls (1.05,1.35) and (1.05,2.25) .. (0,2.25)
   .. controls (-.58,2.25) and (-.58,2.7) .. (0,2.7)
   .. controls (1.3,2.7) and (2.3,2.7) .. (2.3,2.35)
   -- (2.3,0);
 \foreach \j/\y in {1/.45,2/.9,3/1.35,5/2.25,6/2.7}{
  \fill[blue!65!black] (0,\y) circle (2pt);
  \node[anchor=south east,inner sep=2pt] at (0,\y) {$\j$};
 }
 \draw[black!50,fill=white] (0,1.8) circle (2pt);
 \node[anchor=east] at (-.07,1.8) {$4$ vacant};
 \draw[thick,fill=white] (-1.65,0) circle (2pt);
 \draw[thick,fill=white] (2.3,0) circle (2pt);
 \node[below=3pt] at (-1.65,0) {$a$};
 \node[below=3pt] at (2.3,0) {$b$};
\end{tikzpicture}
\caption{A single arch may cross a cut several times. In this schematic
example the left half has a source strand at slot $1$ and pairs $(2,3)$,
$(5,6)$; the right half has pairs $(1,2)$, $(3,5)$ and a source strand
at slot $6$. Slot $4$ is vacant on both sides. Each source strand is
exposed, meaning that no pair encloses its slot. Gluing the two states
recovers one path from $a$ to $b$, with no closed component.}
\label{fig:cut-states}
\end{figure}
